% Self-contained CVPR2022 course milestone; original official style embedded unchanged.
% Style source: https://cvpr2022.thecvf.com/sites/default/files/2021-10/cvpr2022-author_kit-v1_1-1.zip
\begin{filecontents*}{cvpr.sty}
% ---------------------------------------------------------------
%
% The last version of CVPR/ICCV LaTeX template had been developed
% by Paolo.Ienne@di.epfl.ch and awf@acm.org about 15 years ago.
% That version suffered from several issues:
% 1. Authors needed several individual files: cvpr.sty,
%    cvpr_eso.sty, eso-pic.sty.
% 2. For CVPR/ICCV rebuttals, another version of cvpr.sty was
%    required.
% 3. Several warnings arose due to deprecated options.
%
% More recently, Ming-Ming Cheng (cmm_spam@nankai.edu.cn) created
% a single style file that helps to unify review, rebuttal, and
% final versions with a class.
%
% This more recent style has been further modernized for CVPR 2022
% by Stefan Roth (stefan.roth@NOSPAMtu-darmstadt.de).
%
% Acknowledgements: This file is built on the template by
% Ming-Ming Cheng (https://github.com/MCG-NKU/CVPR_Template).
% ---------------------------------------------------------------




% ---------------------------------------------------------------
%
% $Id: cvpr.sty,v 1.3 2005/10/24 19:56:15 awf Exp $
% by Paolo.Ienne@di.epfl.ch some mods by awf@acm.org
%
% ---------------------------------------------------------------
%
% no guarantee is given that the format corresponds perfectly to
% IEEE 8.5" x 11" Proceedings, but most features should be ok.
%
% ---------------------------------------------------------------
% with LaTeX2e:
% =============
%
% use as
%   \documentclass[times,10pt,twocolumn]{article}
%   \usepackage[options]{cvpr}
%   \usepackage{times}
%
% "options" should be replaced by
%  * "review" for submitting a paper for review,
%  * "final" for the camera ready, and
%  * "rebuttal" for the author rebuttal.
%
% specify references as
%   {\small
%   \bibliographystyle{ieee}
%   \bibliography{...your files...}
%   }
% ---------------------------------------------------------------



% ---------------------------------------------------------------
%
%\usepackage{eso-pic}
%
%%
%% This is file `eso-pic.sty',
%% generated with the docstrip utility.
%%
%% The original source files were:
%%
%% eso-pic.dtx  (with options: `package')
%%
%% This is a generated file.
%%
%% Copyright (C) 1998-2002 by Rolf Niepraschk <niepraschk@ptb.de>
%%
%% This file may be distributed and/or modified under the conditions of
%% the LaTeX Project Public License, either version 1.2 of this license
%% or (at your option) any later version.  The latest version of this
%% license is in:
%%
%%    http://www.latex-project.org/lppl.txt
%%
%% and version 1.2 or later is part of all distributions of LaTeX version
%% 1999/12/01 or later.
%%


%
\NeedsTeXFormat{LaTeX2e}[1999/12/01]
\ProvidesPackage{cvpr}[2021/08/23 Example LaTex class for IEEE CVPR]

\RequirePackage{times}    % Integrate Times for here
\RequirePackage{cite}     % Automatically ordered citations
\RequirePackage{xspace}

\RequirePackage{silence}  % Suppress unwanted warnings
\hbadness=10000 \vbadness=10000 \vfuzz=30pt \hfuzz=30pt
\WarningFilter{latexfont}{Font shape declaration}
\WarningFilter{latex}{Font shape}
\WarningFilter{hyperref}{Token not allowed in a PDF string}
\WarningFilter[rebuttal]{latex}{No \author given}
\RequirePackage{etoolbox}

% Use modern caption package to allow for sub-figures etc.
% Reproduces the original CVPR/ICCV style as closely as possible.
\RequirePackage[format=plain,labelformat=simple,labelsep=period,font=small,compatibility=false]{caption}
\RequirePackage[font=footnotesize,skip=3pt,subrefformat=parens]{subcaption}


\newtoggle{cvprfinal}        % Camera-ready version
\newtoggle{cvprrebuttal}     % Rebuttal
\newtoggle{cvprpagenumbers}  % Force page numbers (in camera ready)
\toggletrue{cvprfinal}
\togglefalse{cvprrebuttal}
\togglefalse{cvprpagenumbers}
\DeclareOption{review}{\togglefalse{cvprfinal}\toggletrue{cvprpagenumbers}}
\DeclareOption{rebuttal}{\togglefalse{cvprfinal}\toggletrue{cvprrebuttal}}
\DeclareOption{pagenumbers}{\toggletrue{cvprpagenumbers}}
\DeclareOption*{\PackageWarning{cvpr}{Unkown option `\CurrentOption'}}

\ProcessOptions\relax

% Don't warn about missing author for rebuttal
\iftoggle{cvprrebuttal}{%
  \ActivateWarningFilters[rebuttal]
}{}

% Breaking lines for URLs in the bib
\RequirePackage[hyphens]{url}
\Urlmuskip=0mu plus 1mu\relax


% ---------------------------------------------------------------
%\input{everyshi.sty}
\newcommand{\@EveryShipout@Hook}{}
\newcommand{\@EveryShipout@AtNextHook}{}
\newcommand*{\EveryShipout}[1]
   {\g@addto@macro\@EveryShipout@Hook{#1}}
\newcommand*{\AtNextShipout}[1]
   {\g@addto@macro\@EveryShipout@AtNextHook{#1}}
\newcommand{\@EveryShipout@Shipout}{%
   \afterassignment\@EveryShipout@Test
   \global\setbox\@cclv= %
   }
\newcommand{\@EveryShipout@Test}{%
   \ifvoid\@cclv\relax
      \aftergroup\@EveryShipout@Output
   \else
      \@EveryShipout@Output
   \fi%
   }
\newcommand{\@EveryShipout@Output}{%
   \@EveryShipout@Hook%
   \@EveryShipout@AtNextHook%
      \gdef\@EveryShipout@AtNextHook{}%
   \@EveryShipout@Org@Shipout\box\@cclv%
   }
\newcommand{\@EveryShipout@Org@Shipout}{}
\newcommand*{\@EveryShipout@Init}{%
   \message{ABD: EveryShipout initializing macros}%
   \let\@EveryShipout@Org@Shipout\shipout
   \let\shipout\@EveryShipout@Shipout
   }
\AtBeginDocument{\@EveryShipout@Init}

% ---------------------------------------------------------------



\newcommand\LenToUnit[1]{#1\@gobble}

\newcommand\AtPageUpperLeft[1]{%
  \begingroup
    \@tempdima=0pt\relax\@tempdimb=\ESO@yoffsetI\relax
    \put(\LenToUnit{\@tempdima},\LenToUnit{\@tempdimb}){#1}%
  \endgroup
}
\newcommand\AtPageLowerLeft[1]{\AtPageUpperLeft{%
  \put(0,\LenToUnit{-\paperheight}){#1}}}
\newcommand\AtPageCenter[1]{\AtPageUpperLeft{%
  \put(\LenToUnit{.5\paperwidth},\LenToUnit{-.5\paperheight}){#1}}%
}
\newcommand\AtTextUpperLeft[1]{%
  \begingroup
    \setlength\@tempdima{1in}%
    \ifodd\c@page%
      \advance\@tempdima\oddsidemargin%
    \else%
      \advance\@tempdima\evensidemargin%
    \fi%
    \@tempdimb=\ESO@yoffsetI\relax\advance\@tempdimb-1in\relax%
    \advance\@tempdimb-\topmargin%
    \advance\@tempdimb-\headheight\advance\@tempdimb-\headsep%
    \put(\LenToUnit{\@tempdima},\LenToUnit{\@tempdimb}){#1}%
  \endgroup
}
\newcommand\AtTextLowerLeft[1]{\AtTextUpperLeft{%
  \put(0,\LenToUnit{-\textheight}){#1}}}
\newcommand\AtTextCenter[1]{\AtTextUpperLeft{%
  \put(\LenToUnit{.5\textwidth},\LenToUnit{-.5\textheight}){#1}}}
\newcommand{\ESO@HookI}{} \newcommand{\ESO@HookII}{}
\newcommand{\ESO@HookIII}{}
\newcommand{\AddToShipoutPicture}{%
  \@ifstar{\g@addto@macro\ESO@HookII}{\g@addto@macro\ESO@HookI}}
\newcommand{\ClearShipoutPicture}{\global\let\ESO@HookI\@empty}
\newcommand\ESO@isMEMOIR[1]{}
\@ifclassloaded{memoir}{\renewcommand\ESO@isMEMOIR[1]{#1}}{}
\newcommand{\@ShipoutPicture}{%
  \bgroup
    \@tempswafalse%
    \ifx\ESO@HookI\@empty\else\@tempswatrue\fi%
    \ifx\ESO@HookII\@empty\else\@tempswatrue\fi%
    \ifx\ESO@HookIII\@empty\else\@tempswatrue\fi%
    \if@tempswa%
      \@tempdima=1in\@tempdimb=-\@tempdima%
      \advance\@tempdimb\ESO@yoffsetI%
      \ESO@isMEMOIR{%
        \advance\@tempdima\trimedge%
        \advance\@tempdima\paperwidth%
        \advance\@tempdima-\stockwidth%
        \if@twoside\ifodd\c@page\else%
          \advance\@tempdima-2\trimedge%
          \advance\@tempdima-\paperwidth%
          \advance\@tempdima\stockwidth%
        \fi\fi%
        \advance\@tempdimb\trimtop}%
      \unitlength=1pt%
      \global\setbox\@cclv\vbox{%
        \vbox{\let\protect\relax
          \pictur@(0,0)(\strip@pt\@tempdima,\strip@pt\@tempdimb)%
            \ESO@HookIII\ESO@HookI\ESO@HookII%
            \global\let\ESO@HookII\@empty%
          \endpicture}%
          \nointerlineskip%
        \box\@cclv}%
    \fi
  \egroup
}
\EveryShipout{\@ShipoutPicture}
\RequirePackage{keyval}
\newif\ifESO@dvips\ESO@dvipsfalse \newif\ifESO@grid\ESO@gridfalse
\newif\ifESO@texcoord\ESO@texcoordfalse
\newcommand*\ESO@gridunitname{}
\newcommand*\ESO@gridunit{}
\newcommand*\ESO@labelfactor{}
\newcommand*\ESO@griddelta{}\newcommand*\ESO@griddeltaY{}
\newcommand*\ESO@gridDelta{}\newcommand*\ESO@gridDeltaY{}
\newcommand*\ESO@gridcolor{}
\newcommand*\ESO@subgridcolor{}
\newcommand*\ESO@subgridstyle{dotted}% ???
\newcommand*\ESO@gap{}
\newcommand*\ESO@yoffsetI{}\newcommand*\ESO@yoffsetII{}
\newcommand*\ESO@gridlines{\thinlines}
\newcommand*\ESO@subgridlines{\thinlines}
\newcommand*\ESO@hline[1]{\ESO@subgridlines\line(1,0){#1}}
\newcommand*\ESO@vline[1]{\ESO@subgridlines\line(0,1){#1}}
\newcommand*\ESO@Hline[1]{\ESO@gridlines\line(1,0){#1}}
\newcommand*\ESO@Vline[1]{\ESO@gridlines\line(0,1){#1}}
\newcommand\ESO@fcolorbox[4][]{\fbox{#4}}
\newcommand\ESO@color[1]{}
\newcommand\ESO@colorbox[3][]{%
  \begingroup
    \fboxrule=0pt\fbox{#3}%
  \endgroup
}
\newcommand\gridSetup[6][]{%
  \edef\ESO@gridunitname{#1}\edef\ESO@gridunit{#2}
  \edef\ESO@labelfactor{#3}\edef\ESO@griddelta{#4}
  \edef\ESO@gridDelta{#5}\edef\ESO@gap{#6}}
\define@key{ESO}{texcoord}[true]{\csname ESO@texcoord#1\endcsname}
\define@key{ESO}{pscoord}[true]{\csname @tempswa#1\endcsname
  \if@tempswa\ESO@texcoordfalse\else\ESO@texcoordtrue\fi}
\define@key{ESO}{dvips}[true]{\csname ESO@dvips#1\endcsname}
\define@key{ESO}{grid}[true]{\csname ESO@grid#1\endcsname
  \setkeys{ESO}{gridcolor=black,subgridcolor=black}}
\define@key{ESO}{colorgrid}[true]{\csname ESO@grid#1\endcsname
  \setkeys{ESO}{gridcolor=red,subgridcolor=green}}
\define@key{ESO}{gridcolor}{\def\ESO@gridcolor{#1}}
\define@key{ESO}{subgridcolor}{\def\ESO@subgridcolor{#1}}
\define@key{ESO}{subgridstyle}{\def\ESO@subgridstyle{#1}}%
\define@key{ESO}{gridunit}{%
  \def\@tempa{#1}
  \def\@tempb{bp}
  \ifx\@tempa\@tempb
    \gridSetup[\@tempa]{1bp}{1}{10}{50}{2}
  \else
    \def\@tempb{pt}
    \ifx\@tempa\@tempb
      \gridSetup[\@tempa]{1pt}{1}{10}{50}{2}
    \else
      \def\@tempb{in}
      \ifx\@tempa\@tempb
        \gridSetup[\@tempa]{.1in}{.1}{2}{10}{.5}
      \else
        \gridSetup[mm]{1mm}{1}{5}{20}{1}
      \fi
    \fi
  \fi
}
%\setkeys{ESO}{subgridstyle=solid,pscoord=true,gridunit=mm}
% \def\ProcessOptionsWithKV#1{%
%   \let\@tempc\@empty
%   \@for\CurrentOption:=\@classoptionslist\do{%
%     \@ifundefined{KV@#1@\CurrentOption}%
%     {}{\edef\@tempc{\@tempc,\CurrentOption,}}}%
%   \edef\@tempc{%
%     \noexpand\setkeys{#1}{\@tempc\@ptionlist{\@currname.\@currext}}
%   }%
%   \@tempc
%   \AtEndOfPackage{\let\@unprocessedoptions\relax}}%
%\ProcessOptionsWithKV{ESO}%
\newcommand\ESO@div[2]{%
  \@tempdima=#1\relax\@tempdimb=\ESO@gridunit\relax
  \@tempdimb=#2\@tempdimb\divide\@tempdima by \@tempdimb%
  \@tempcnta\@tempdima\advance\@tempcnta\@ne}
\AtBeginDocument{%
  \IfFileExists{color.sty}
  {%
    \RequirePackage{color}
    \let\ESO@color=\color\let\ESO@colorbox=\colorbox
    \let\ESO@fcolorbox=\fcolorbox
  }{}
  \@ifundefined{Gin@driver}{}%
  {%
    \ifx\Gin@driver\@empty\else%
      \filename@parse{\Gin@driver}\def\reserved@a{dvips}%
      \ifx\filename@base\reserved@a\ESO@dvipstrue\fi%
    \fi
  }%
  \ifx\pdfoutput\undefined\else
    \ifx\pdfoutput\relax\else
      \ifcase\pdfoutput\else
        \ESO@dvipsfalse%
      \fi
    \fi
  \fi
  \ifESO@dvips\def\@tempb{eepic}\else\def\@tempb{epic}\fi
  \def\@tempa{dotted}%\def\ESO@gap{\LenToUnit{6\@wholewidth}}%
  \ifx\@tempa\ESO@subgridstyle
    \IfFileExists{\@tempb.sty}%
    {%
      \RequirePackage{\@tempb}
      \renewcommand*\ESO@hline[1]{\ESO@subgridlines\dottedline{\ESO@gap}%
        (0,0)(##1,0)}
      \renewcommand*\ESO@vline[1]{\ESO@subgridlines\dottedline{\ESO@gap}%
        (0,0)(0,##1)}
    }{}
  \else
    \ifx\ESO@gridcolor\ESO@subgridcolor%
      \renewcommand*\ESO@gridlines{\thicklines}
    \fi
  \fi
}
\ifESO@texcoord
  \def\ESO@yoffsetI{0pt}\def\ESO@yoffsetII{-\paperheight}
  \edef\ESO@griddeltaY{-\ESO@griddelta}\edef\ESO@gridDeltaY{-\ESO@gridDelta}
\else
  \def\ESO@yoffsetI{\paperheight}\def\ESO@yoffsetII{0pt}
  \edef\ESO@griddeltaY{\ESO@griddelta}\edef\ESO@gridDeltaY{\ESO@gridDelta}
\fi
\newcommand\ESO@gridpicture{%
  \begingroup
    \setlength\unitlength{\ESO@gridunit}%
    \ESO@color{\ESO@subgridcolor}%
    \ESO@div{\paperheight}{\ESO@griddelta}%
    \multiput(0,0)(0,\ESO@griddeltaY){\@tempcnta}%
      {\ESO@hline{\LenToUnit{\paperwidth}}}%
    \ESO@div{\paperwidth}{\ESO@griddelta}%
    \multiput(0,\LenToUnit{\ESO@yoffsetII})(\ESO@griddelta,0){\@tempcnta}%
      {\ESO@vline{\LenToUnit{\paperheight}}}%
    \ESO@color{\ESO@gridcolor}%
    \ESO@div{\paperheight}{\ESO@gridDelta}%
    \multiput(0,0)(0,\ESO@gridDeltaY){\@tempcnta}%
      {\ESO@Hline{\LenToUnit{\paperwidth}}}%
    \ESO@div{\paperwidth}{\ESO@gridDelta}%
    \multiput(0,\LenToUnit{\ESO@yoffsetII})(\ESO@gridDelta,0){\@tempcnta}%
      {\ESO@Vline{\LenToUnit{\paperheight}}}%
    \fontsize{10}{12}\normalfont%
    \ESO@div{\paperwidth}{\ESO@gridDelta}%
    \multiput(0,\ESO@gridDeltaY)(\ESO@gridDelta,0){\@tempcnta}{%
      \@tempcntb=\@tempcnta\advance\@tempcntb-\@multicnt%
      \ifnum\@tempcntb>1\relax
        \multiply\@tempcntb by \ESO@gridDelta\relax%
        \@tempdima=\@tempcntb sp\@tempdima=\ESO@labelfactor\@tempdima%
        \@tempcntb=\@tempdima%
        \makebox(0,0)[c]{\ESO@colorbox{white}{\the\@tempcntb}}%
      \fi}%
    \ifx\ESO@gridunitname\@empty\def\@tempa{0}\else\def\@tempa{1}\fi%
    \ESO@div{\paperheight}{\ESO@gridDelta}%
    \multiput(\ESO@gridDelta,0)(0,\ESO@gridDeltaY){\@tempcnta}{%
      \@tempcntb=\@tempcnta\advance\@tempcntb-\@multicnt%
      \ifnum\@tempcntb>\@tempa\relax
        \multiply\@tempcntb by \ESO@gridDelta\relax%
        \@tempdima=\@tempcntb sp\@tempdima=\ESO@labelfactor\@tempdima%
        \@tempcntb=\@tempdima%
        \makebox(0,0)[c]{\ESO@colorbox{white}{\the\@tempcntb}}%
      \fi
    }%
    \ifx\ESO@gridunitname\@empty\else%
      \thicklines\fboxrule=\@wholewidth%
      \put(\ESO@gridDelta,\ESO@gridDeltaY){\makebox(0,0)[c]{%
        \ESO@fcolorbox{\ESO@gridcolor}{white}{%
          \textbf{\ESO@gridunitname}}}}%
    \fi
    \normalcolor%
  \endgroup
}
\ifESO@grid\g@addto@macro\ESO@HookIII{\ESO@gridpicture}\fi
% ---------------------------------------------------------------




\typeout{CVPR 8.5 x 11-Inch Proceedings Style `cvpr.sty'.}

% ten point helvetica bold required for captions
% eleven point times bold required for second-order headings
% in some sites the name of the fonts may differ,
% change the name here:
\font\cvprtenhv  = phvb at 8pt % *** IF THIS FAILS, SEE cvpr.sty ***
\font\elvbf  = ptmb scaled 1100

% If the above lines give an error message, try to comment them and
% uncomment these:
%\font\cvprtenhv  = phvb7t at 8pt
%\font\elvbf  = ptmb7t scaled 1100

% set dimensions of columns, gap between columns, and paragraph indent
\setlength{\textheight}{8.875in}
\setlength{\textwidth}{6.875in}
\setlength{\columnsep}{0.3125in}
\setlength{\topmargin}{0in}
\setlength{\headheight}{0in}
\setlength{\headsep}{0in}
\setlength{\parindent}{1pc}
\setlength{\oddsidemargin}{-.304in}
\setlength{\evensidemargin}{-.304in}


% memento from size10.clo
% \normalsize{\@setfontsize\normalsize\@xpt\@xiipt}
% \small{\@setfontsize\small\@ixpt{11}}
% \footnotesize{\@setfontsize\footnotesize\@viiipt{9.5}}
% \scriptsize{\@setfontsize\scriptsize\@viipt\@viiipt}
% \tiny{\@setfontsize\tiny\@vpt\@vipt}
% \large{\@setfontsize\large\@xiipt{14}}
% \Large{\@setfontsize\Large\@xivpt{18}}
% \LARGE{\@setfontsize\LARGE\@xviipt{22}}
% \huge{\@setfontsize\huge\@xxpt{25}}
% \Huge{\@setfontsize\Huge\@xxvpt{30}}


% Suppress page numbers when the appropriate option is given
\iftoggle{cvprpagenumbers}{}{%
  \pagestyle{empty}
}

\AtBeginDocument{%
  % Print an error if document class other than article is used
  \@ifclassloaded{article}{}{%
    \PackageError{cvpr}{Package only meant to be used with document class `article'}{Change document class to `article'.}
  }
  % Print a warning if incorrect options for article are specified
  \@ifclasswith{article}{10pt}{}{%
    \PackageWarningNoLine{cvpr}{Incorrect font size specified - CVPR requires 10-point fonts. Please load document class `article' with `10pt' option}
  }
  \@ifclasswith{article}{twocolumn}{}{%
    \PackageWarningNoLine{cvpr}{Single column document - CVPR requires papers to have two-column layout. Please load document class `article' with `twocolumn' option}
  }
  \@ifclasswith{article}{letterpaper}{}{%
    \PackageWarningNoLine{cvpr}{Incorrect paper size - CVPR uses paper size `letter'. Please load document class `article' with `letterpaper' option}
  }
  % Print a warning if hyperref is not loaded and/or if the pagebackref option is missing
  \iftoggle{cvprfinal}{%
    \@ifpackageloaded{hyperref}{}{%
      \PackageWarningNoLine{cvpr}{Package `hyperref' is not loaded, but highly recommended for camera-ready version}
    }
  }{%
    \@ifpackageloaded{hyperref}{
      \@ifpackagewith{hyperref}{pagebackref}{}{
        \PackageWarningNoLine{cvpr}{Package `hyperref' is not loaded with option `pagebackref', which is strongly recommended for review version}
      }
    }{%
      \PackageWarningNoLine{cvpr}{Package `hyperref' is not loaded, but strongly recommended for review version}
    }
  }
}

\def\@maketitle
   {
   \newpage
   \null
   \iftoggle{cvprrebuttal}{\vspace*{-.3in}}{\vskip .375in}
   \begin{center}
      % smaller title font only for rebuttal
      \iftoggle{cvprrebuttal}{{\large \bf \@title \par}}{{\Large \bf \@title \par}}
      % additional two empty lines at the end of the title
      \iftoggle{cvprrebuttal}{\vspace*{-22pt}}{\vspace*{24pt}}
      {
      \large
      \lineskip .5em
      \begin{tabular}[t]{c}
        \iftoggle{cvprfinal}{
          \@author
        }{
          \iftoggle{cvprrebuttal}{}{
            Anonymous \confName~submission\\
            \vspace*{1pt}\\
            Paper ID \cvprPaperID
          }
        }
      \end{tabular}
      \par
      }
      % additional small space at the end of the author name
      \vskip .5em
      % additional empty line at the end of the title block
      \vspace*{12pt}
   \end{center}
   }

\def\abstract
   {%
   % Suppress page numbers when the appropriate option is given
   \iftoggle{cvprpagenumbers}{}{%
     \thispagestyle{empty}
   }
   \centerline{\large\bf Abstract}%
   \vspace*{12pt}%
   \it%
   }

\def\endabstract
   {
   % additional empty line at the end of the abstract
   \vspace*{12pt}
   }

\def\affiliation#1{\gdef\@affiliation{#1}} \gdef\@affiliation{}

% correct heading spacing and type
\def\cvprsection{\@startsection {section}{1}{\z@}
   {10pt plus 2pt minus 2pt}{7pt} {\large\bf}}
\def\cvprssect#1{\cvprsection*{#1}}
\def\cvprsect#1{\cvprsection{\hskip -1em.~#1}}
\def\section{\@ifstar\cvprssect\cvprsect}

\def\cvprsubsection{\@startsection {subsection}{2}{\z@}
   {8pt plus 2pt minus 2pt}{6pt} {\elvbf}}
\def\cvprssubsect#1{\cvprsubsection*{#1}}
\def\cvprsubsect#1{\cvprsubsection{\hskip -1em.~#1}}
\def\subsection{\@ifstar\cvprssubsect\cvprsubsect}

%% --------- Page background marks: Ruler and confidentiality

% ----- define vruler
\makeatletter
\newbox\cvprrulerbox
\newcount\cvprrulercount
\newdimen\cvprruleroffset
\newdimen\cv@lineheight
\newdimen\cv@boxheight
\newbox\cv@tmpbox
\newcount\cv@refno
\newcount\cv@tot
% NUMBER with left flushed zeros  \fillzeros[<WIDTH>]<NUMBER>
\newcount\cv@tmpc@ \newcount\cv@tmpc
\def\fillzeros[#1]#2{\cv@tmpc@=#2\relax\ifnum\cv@tmpc@<0\cv@tmpc@=-\cv@tmpc@\fi
\cv@tmpc=1 %
\loop\ifnum\cv@tmpc@<10 \else \divide\cv@tmpc@ by 10 \advance\cv@tmpc by 1 \fi
   \ifnum\cv@tmpc@=10\relax\cv@tmpc@=11\relax\fi \ifnum\cv@tmpc@>10 \repeat
\ifnum#2<0\advance\cv@tmpc1\relax-\fi
\loop\ifnum\cv@tmpc<#1\relax0\advance\cv@tmpc1\relax\fi \ifnum\cv@tmpc<#1 \repeat
\cv@tmpc@=#2\relax\ifnum\cv@tmpc@<0\cv@tmpc@=-\cv@tmpc@\fi \relax\the\cv@tmpc@}%
% \makevruler[<SCALE>][<INITIAL_COUNT>][<STEP>][<DIGITS>][<HEIGHT>]
\def\makevruler[#1][#2][#3][#4][#5]{\begingroup\offinterlineskip
\textheight=#5\vbadness=10000\vfuzz=120ex\overfullrule=0pt%
\global\setbox\cvprrulerbox=\vbox to \textheight{%
{\parskip=0pt\hfuzz=150em\cv@boxheight=\textheight
\cv@lineheight=#1\global\cvprrulercount=#2%
\cv@tot\cv@boxheight\divide\cv@tot\cv@lineheight\advance\cv@tot2%
\cv@refno1\vskip-\cv@lineheight\vskip1ex%
\loop\setbox\cv@tmpbox=\hbox to0cm{{\cvprtenhv\hfil\fillzeros[#4]\cvprrulercount}}%
\ht\cv@tmpbox\cv@lineheight\dp\cv@tmpbox0pt\box\cv@tmpbox\break
\advance\cv@refno1\global\advance\cvprrulercount#3\relax
\ifnum\cv@refno<\cv@tot\repeat}}\endgroup}%
\makeatother
% ----- end of vruler

% \makevruler[<SCALE>][<INITIAL_COUNT>][<STEP>][<DIGITS>][<HEIGHT>]
\def\cvprruler#1{\makevruler[12pt][#1][1][3][0.993\textheight]\usebox{\cvprrulerbox}}
\AddToShipoutPicture{%
  \iftoggle{cvprfinal}{
  }{
    \cvprruleroffset=\textheight
    \advance\cvprruleroffset by -3.7pt
      \color[rgb]{.5,.5,1}
      \AtTextUpperLeft{%
        \put(\LenToUnit{-35pt},\LenToUnit{-\cvprruleroffset}){%left ruler
          \cvprruler{\cvprrulercount}}
        %\put(\LenToUnit{\textwidth\kern 30pt},\LenToUnit{-\cvprruleroffset})
        \put(\LenToUnit{\dimexpr \textwidth+30pt},\LenToUnit{-\cvprruleroffset}){%right ruler
          \cvprruler{\cvprrulercount}}
      }
    \def\pid{\parbox{1in}{\begin{center}\bf\sf{\small \confName}\\\#\cvprPaperID\end{center}}}
      \AtTextUpperLeft{%paperID in corners
        \put(\LenToUnit{-65pt},\LenToUnit{45pt}){\pid}
        \put(\LenToUnit{\textwidth\kern-8pt},\LenToUnit{45pt}){\pid}
      }
      \AtTextUpperLeft{%confidential
        \put(0,\LenToUnit{1cm}){\parbox{\textwidth}{\centering\cvprtenhv
        \confName~\confYear~Submission \#\cvprPaperID. CONFIDENTIAL REVIEW COPY.  DO NOT DISTRIBUTE.}}
      }
  }
}

%%% Make figure placement a little more predictable.
% We trust the user to move figures if this results
% in ugliness.
% Minimize bad page breaks at figures
\renewcommand{\textfraction}{0.01}
\renewcommand{\floatpagefraction}{0.99}
\renewcommand{\topfraction}{0.99}
\renewcommand{\bottomfraction}{0.99}
\renewcommand{\dblfloatpagefraction}{0.99}
\renewcommand{\dbltopfraction}{0.99}
\setcounter{totalnumber}{99}
\setcounter{topnumber}{99}
\setcounter{bottomnumber}{99}

% Add a period to the end of an abbreviation unless there's one
% already, then \xspace.
\makeatletter
\DeclareRobustCommand\onedot{\futurelet\@let@token\@onedot}
\def\@onedot{\ifx\@let@token.\else.\null\fi\xspace}

\def\eg{\emph{e.g}\onedot} \def\Eg{\emph{E.g}\onedot}
\def\ie{\emph{i.e}\onedot} \def\Ie{\emph{I.e}\onedot}
\def\cf{\emph{cf}\onedot} \def\Cf{\emph{Cf}\onedot}
\def\etc{\emph{etc}\onedot} \def\vs{\emph{vs}\onedot}
\def\wrt{w.r.t\onedot} \def\dof{d.o.f\onedot}
\def\iid{i.i.d\onedot} \def\wolog{w.l.o.g\onedot}
\def\etal{\emph{et al}\onedot}
\makeatother

% ---------------------------------------------------------------

\end{filecontents*}
\documentclass[10pt,twocolumn,letterpaper]{article}
\usepackage[T1]{fontenc}
\usepackage[pagenumbers]{cvpr}
\usepackage{amsmath,amssymb,booktabs}
\usepackage{pgfplots}
\pgfplotsset{compat=1.18}
\usepackage[breaklinks=true,colorlinks=true,citecolor=blue,urlcolor=blue,linkcolor=blue]{hyperref}
\def\cvprPaperID{}
\def\confName{ELEC4240}
\def\confYear{2026}
\title{Low-Cost Mixed-Resolution Adaptation of Marigold\\for Indoor Depth Estimation}
\author{HO, Chun Wai\\21053878\and Tong, Man Hung\\21064669\and MA, Shenhan\\21041382}
\begin{document}
\maketitle

\section{Introduction}
Monocular depth estimation predicts scene geometry from a single RGB image. It is useful for indoor perception, but collecting depth supervision and adapting large models can be expensive. Marigold repurposes a diffusion image generator for relative depth estimation~\cite{marigold}. Our project asks whether a small, inexpensive adaptation of an existing depth model can improve indoor predictions on an 8\,GB laptop GPU, and whether any improvement transfers to unseen scenes.

We start from Marigold Depth v1.1~\cite{marigold25} rather than training a generator from scratch. We implement low-rank adaptation (LoRA)~\cite{lora} and compare fixed high-resolution training with alternating low/high-resolution updates. The study combines controlled training repetitions, held-out evaluation and an external validation cohort. Our contribution is this adaptation and evaluation study; we do not claim to invent LoRA, Marigold or a new diffusion architecture.

At this milestone, the training and evaluation pipeline is operational. On 200 external scene groups, mixed-resolution training reduces aligned absolute relative error (AbsRel) by 13.3\% at 256-pixel inference, with support after correction for multiple comparisons. At 512 pixels, the mean gain is smaller and remains statistically unconfirmed. These outcomes motivate studying when adaptation helps and why relative depth improvements do not necessarily yield accurate distances.

\section{Problem Statement}
\paragraph{Task and data.}
Given an RGB image $x$, predict a dense depth field $f_\theta(x)$. We use labeled indoor RGB-D images from NYU Depth v2~\cite{nyu} for adaptation. The latest replication uses three random draws of 128 training scenes from 217 eligible scenes, with one selected frame per scene. The draws can overlap; they are not independent domains. A fixed 32-scene validation set supplies calibration parameters. A 31-scene cohort, unused in earlier project phases when the replication began, is the primary NYUv2 test; a previously examined 64-scene cohort is kept separate.

For external validation, we select 200 distinct source-space groups from the SUN3D branch of SUN RGB-D~\cite{sunrgbd,sun3d}, one frame per group. NYUv2-containing and other source branches are excluded. Groups and candidate frames are ranked by fixed hashes before prediction. Raw measured-depth pairs must decode correctly, share image dimensions, and contain enough valid pixels. Exact RGB duplicates of earlier local inputs are excluded. The selected groups map to 63 broader location proxies, which are used to assess dependence between nearby scenes. This is a custom external subset, not the full SUN RGB-D benchmark.

\paragraph{Expected outcomes and evaluation.}
We hypothesize that mixed-resolution updates improve relative depth while retaining a small adaptation budget. An improvement is an experimental question, not a prerequisite for success. The primary metric is mean per-image AbsRel:
\begin{equation}
\mathrm{AbsRel}=\frac{1}{|V|}\sum_{i\in V}\frac{|\hat d_i-d_i|}{d_i},
\end{equation}
where $V$ contains valid target pixels. Because Marigold predicts relative depth, we first fit a least-squares scale and shift to each image's ground-truth (GT) depth for the aligned evaluation. We also report RMSE and the fraction satisfying $\max(\hat d_i/d_i,d_i/\hat d_i)<1.25$, denoted $\delta_1$.

Aligned scores measure relative structure, not deployable metric accuracy. A separate diagnostic uses one affine calibration per model and resolution, learned only on the 32 NYUv2 validation scenes and then frozen. External evaluation uses all measured pixels with $0.1<d<10$\,m; the earlier NYUv2 evaluation uses its fixed crop. These protocols are not pooled. Qualitative depth/error maps and measured time and peak allocated GPU memory complement the quantitative analysis.

\section{Technical Approach}
\paragraph{Implementation and training objective.}
Our compact trainer follows the public Marigold latent-denoising formulation and uses its pretrained weights through Diffusers; adapters use PEFT.\footnote{Base implementations: \url{https://github.com/prs-eth/Marigold}, \url{https://github.com/huggingface/diffusers}, and \url{https://github.com/huggingface/peft}.} The VAE, text encoder and original U-Net weights remain frozen. Rank-4 LoRA with scaling parameter 4 updates the U-Net attention query, key, value and output projections, totaling 829,952 trainable parameters.

RGB and depth are encoded into cached VAE latents. Depth is normalized per image using its 2nd and 98th percentiles. At diffusion timestep $t$, noisy depth latents are concatenated with RGB latents. The adapter minimizes masked velocity-prediction error:
\begin{equation}
\mathcal L=\frac{1}{|M|}\sum_{j\in M}\left\|v_{\theta,\phi}(z_x,z_{d,t},t)_j-v_{t,j}\right\|^2,
\end{equation}
where $\theta$ is frozen, $\phi$ contains LoRA parameters, and $M$ excludes latent cells affected by invalid depth. This is an existing denoising objective, not a new geometric consistency loss.

\paragraph{Controlled adaptation variants.}
\textit{High-only} uses a 512-pixel processing long side for every update; \textit{mixed} alternates 256 and 512, giving 160 updates at each resolution. Both use 320 total updates, batch size 1, AdamW with learning rate $10^{-4}$ and weight decay 0.01, BF16 autocasting and gradient checkpointing. We pair image order and timestep schedules across methods. Equal update counts do not imply equal compute. Five seeds per training draw yield $3\times5\times2=30$ adapters. Inference uses one denoising step, ensemble size 1 and fixed per-image random seeds at both resolutions.

Original Marigold is the main reference. Depth Anything V2's indoor Small metric model~\cite{dav2}, fine-tuned on Hypersim, is an additional descriptive reference. Its official default preprocessing uses a larger pixel budget and different prior supervision; it is not a matched-cost competitor in the formal comparison family.

\paragraph{Frozen validation and uncertainty.}
Before external inference, we froze all 30 checkpoints, scene IDs, metrics and the stopping rule. The 63 conditions comprise original Marigold and 30 adapters at two resolutions, plus the specialist, giving 12,600 predictions. Six comparisons cover mixed versus original, high-only versus original, and mixed versus high-only at each resolution. Each uses centered two-sided tests with 20,000 paired bootstrap replicates with training-draw, seed and scene resampling. We apply Holm correction~\cite{holm} separately to nested and crossed resampling analyses. A directional improvement requires a negative mean error difference and adjusted $p<0.05$ in both. A third analysis resamples broader location groups. With only three training draws and imperfect location metadata, uncertainty estimates remain approximate.

\section{Intermediate Results}
\paragraph{NYUv2 replication.}
On the fresh 31-scene cohort, mixed training at 512 inference reduces aligned AbsRel from 0.09687 to 0.08462 (12.6\%). However, nested/crossed Holm-adjusted $p=0.1088/0.1094$ does not confirm this improvement. At 256 inference, mixed training improves over high-only training (0.10128 versus 0.12696; adjusted $p=0.0042/0.0084$). These results led to a prospective external test rather than reusing previously examined test images as new evidence.

\begin{table}[t]
\centering\small
\caption{External cohort: aligned scores averaged over 200 scene groups. Adapter scores average all 15 draw/seed runs. H: high-only training; M: mixed training.}
\label{tab:scores}
\begin{tabular}{lrrr}\toprule
Method / inference & AbsRel $\downarrow$ & RMSE $\downarrow$ & $\delta_1\uparrow$\\\midrule
Original / 256 & .11106 & .35647 & .88161\\
H / 256 & .11188 & .35673 & .87626\\
M / 256 & .09632 & .31430 & .90983\\\midrule
Original / 512 & .07224 & .25609 & .94551\\
H / 512 & .07154 & .24980 & .94506\\
M / 512 & .07038 & .24634 & .94669\\\midrule
Specialist / default & .07528 & .26772 & .93997\\\bottomrule
\end{tabular}
\end{table}

\paragraph{External generalization.}
Table~\ref{tab:scores} and Fig.~\ref{fig:gain} summarize the fixed external cohort. At 256 inference, mixed training improves over original Marigold by 13.3\% (nested/crossed Holm $p=0.00030/0.00030$); the location sensitivity also supports improvement ($p=0.00150$). Mixed also improves over high-only training at 256 (all three adjusted $p=0.00030$). At 512, mixed improves the mean by only 2.6\%, with nested/crossed/location adjusted $p=0.29684/0.30778/0.61617$. None of the three 512 comparisons passes the predefined rule. High-only versus original at 256 is also unconfirmed ($p=0.99245/1.00000$). Table~\ref{tab:holm} reports the entire comparison family.

\begin{figure}[t]
\centering
\begin{tikzpicture}
\begin{axis}[width=.90\columnwidth,height=4.7cm,scaled x ticks=false,xmin=-.027,xmax=.009,
 ymin=.5,ymax=2.5,ytick={1,2},yticklabels={512 inference,256 inference},
 xlabel={Mixed $-$ original aligned AbsRel},xtick={-.02,-.01,0},
 xticklabels={$-0.02$,$-0.01$,$0$},
 xmajorgrids,grid style={gray!20},axis lines=left,
 tick label style={font=\small},label style={font=\small}]
\addplot[black,dashed] coordinates {(0,.5)(0,2.5)};
\addplot[blue!70!black,very thick] coordinates {(-.02002,2)(-.00970,2)};
\addplot[orange!85!black,very thick] coordinates {(-.00408,1)(.00032,1)};
\addplot[only marks,mark=*,blue!70!black] coordinates {(-.01474,2)};
\addplot[only marks,mark=*,orange!85!black] coordinates {(-.00186,1)};
\end{axis}
\end{tikzpicture}
\caption{External paired error differences with nested-bootstrap percentile 95\% intervals; negative favors mixed training. These are descriptive, unadjusted intervals. Decisions use the complete six-comparison Holm family, not interval exclusion alone.}
\label{fig:gain}
\end{figure}


\begin{table}[t]
\centering\small
\caption{All six external contrasts: Holm-adjusted $p$-values for nested ($p_N$), crossed ($p_C$), and location ($p_L$) resampling. O: original; H/M as in Table~\ref{tab:scores}. Only the two 256-pixel mixed contrasts support improvement.}
\label{tab:holm}
\begin{tabular}{lrrr}\toprule
Contrast (inference) & $p_N$ & $p_C$ & $p_L$\\\midrule
M $-$ O (256) & .00030 & .00030 & .00150\\
M $-$ O (512) & .29684 & .30778 & .61617\\
H $-$ O (256) & .99245 & 1.00000 & 1.00000\\
H $-$ O (512) & .99245 & 1.00000 & 1.00000\\
M $-$ H (256) & .00030 & .00030 & .00030\\
M $-$ H (512) & .26639 & .28039 & .29439\\\bottomrule
\end{tabular}
\end{table}

\paragraph{Why more scenes did not guarantee significance.}
A planning simulation using earlier NYUv2 data estimated 84.3\% power at 200 groups under an assumed AbsRel gain of 0.009. Inflating error SD by 1.5 reduced that estimate to 38.5\%. These are conditional pilot-distribution approximations, not guaranteed power under domain shift. The observed external mixed512 gain is only 0.00186, much smaller than assumed. We therefore retain the nonsignificant result and do not extend this test until it passes.

\paragraph{Metric depth and qualitative observations.}
With fixed NYUv2 calibration, external mixed512 AbsRel is 0.37886 and RMSE is 0.94306\,m; the specialist obtains 0.18231 and 0.47903\,m. Its native, uncalibrated AbsRel is 0.29033. Thus favorable GT-aligned structure does not establish accurate absolute distances. The first three external manifest images and one fixed adapter were selected for visualization before outcomes were inspected. Their depth and error maps show plausible relative structure alongside substantial calibration errors; they are illustrative examples, not a representative failure-rate estimate.

\paragraph{Feasibility and verification.}
The 30 replication trainings used 52.40 minutes of optimizer time and peaked at 2,512\,MiB allocated training memory on an RTX 5060 Laptop GPU. The external study reused these adapters with no additional training; measured inference totaled 32.10 minutes and peaked at 2,946.7\,MiB. Loading, latent caching and audits are excluded from these times. All 12,600 external predictions were checked and their metrics recomputed; 200 raw RGB/depth pairs were independently redecoded. Data, checkpoint and code hashes support reproducibility. These are local measurements, not a controlled hardware benchmark.

\section{Limitations and Remaining Work}
The current evidence supports a low-resolution benefit on this cohort, not universal superiority. A significant result at 256 and a nonsignificant result at 512 do not by themselves establish a resolution interaction. Scene groups may share physical locations, training draws overlap, and broad pretraining overlap cannot be excluded. The study also uses a small adaptation budget and custom evaluation subsets.

For the final report, we plan to (1) test a paired resolution-by-method interaction as a clearly labeled exploratory analysis; (2) study absolute-depth calibration using training/validation data only, including depth-range and invalid-pixel failure analysis; and (3) compare update-matched and compute-matched adaptation budgets. Any new method selection will use development data. Claims from the completed external test will remain fixed; a new confirmatory claim would require a separately frozen, unseen cohort. The final write-up will consolidate ablations and qualitative failures rather than present nonsignificance as equivalence.

{\small
\begin{thebibliography}{9}
\bibitem{marigold} Bingxin Ke et al. Repurposing diffusion-based image generators for monocular depth estimation. In \emph{CVPR}, 2024.
\bibitem{marigold25} Bingxin Ke et al. Marigold: Affordable adaptation of diffusion-based image generators for image analysis. arXiv:2505.09358, 2025.
\bibitem{lora} Edward J. Hu et al. LoRA: Low-rank adaptation of large language models. In \emph{ICLR}, 2022.
\bibitem{nyu} Nathan Silberman, Derek Hoiem, Pushmeet Kohli, and Rob Fergus. Indoor segmentation and support inference from RGBD images. In \emph{ECCV}, 2012.
\bibitem{sunrgbd} Shuran Song, Samuel P. Lichtenberg, and Jianxiong Xiao. SUN RGB-D: A RGB-D scene understanding benchmark suite. In \emph{CVPR}, 2015.
\bibitem{sun3d} Jianxiong Xiao, Andrew Owens, and Antonio Torralba. SUN3D: A database of big spaces reconstructed using SfM and object labels. In \emph{ICCV}, 2013.
\bibitem{dav2} Lihe Yang et al. Depth Anything V2. In \emph{NeurIPS}, 2024.
\bibitem{holm} Sture Holm. A simple sequentially rejective multiple test procedure. \emph{Scandinavian Journal of Statistics}, 6(2):65--70, 1979.
\end{thebibliography}}
\end{document}
